
% Explain Why did this happen? this indicates...
% Key Points: 
% \begin{itemize}
%     \item Key Results from the nominal model, errors and compute times why are they reasonable and acceptable.
%     \item Why we should use combined $\Delta x$ and $\Delta t$ refinement and how that limits credibility 
%     \item Why the simplified geometries for CV and SV are simple but still complex enough
%     \item Reproducibility challenges with verification: large datasets (meshes & outputs) significant computational resources to run the finer meshes  
% \end{itemize}

This study presents a rigorous verification process for a conjugate heat transfer model of a simplified medical device geometry. The objective was to establish confidence in the numerical implementation and to quantify sensitivity to spatial and temporal discretization, as well as solver convergence. These verification activities form a critical step toward building model credibility prior to validation.

Verification plays a particularly important role in medical device modeling, where experimental validation is often limited by cost, feasibility, or measurement uncertainty. In such cases, verification provides a quantitative framework to isolate numerical error, ensuring that remaining uncertainty can be attributed to physical modeling assumptions or experimental limitations rather than deficiencies in the numerical solution.

Code verification results confirm that the numerical implementation accurately represents the governing equations. The hydrodynamic case exhibited an observed order of convergence of $p = 1.90$, with errors decreasing from $\mathcal{O}(10^{-3})$ to $\mathcal{O}(10^{-5})$, while the radial conduction case achieved $p = 2.08$, consistent with second-order spatial accuracy. These results demonstrate that both the momentum and energy equations are correctly discretized and solved under idealized conditions.

The verification process was intentionally structured to progressively isolate the governing physics relevant to the final coupled transient configuration. The hydrodynamic verification case isolated the momentum equations and confirmed accurate velocity-field discretization independent of thermal effects. The radial conduction case then isolated steady-state thermal transport within the solid domain by removing fluid effects through an adiabatic interface boundary condition. While these studies independently verified the momentum and thermal governing equations, neither configuration evaluated transient thermal transport or the thermal coupling behavior at the fluid-solid interface. To bridge this gap, a transient thermal diffusion verification case was introduced by imposing zero velocity throughout the fluid region, reducing the governing energy equation to a pure diffusion problem with an analytical Bessel-function solution. This intermediate verification configuration enabled assessment of transient thermal transport behavior independently of advection while also verifying that the transient solid temperature response was correctly imposed across the fluid-solid interface into the fluid region. The transient diffusion study demonstrated near second-order spatial convergence, with observed orders ranging from $p=2.15$ to $p=2.47$ across multiple transient times.

Solution verification results highlight the importance of evaluating discretization error in coupled, application-relevant configurations. In the transient simplified EDDS case, co-refined mesh and timestep studies revealed that spatial and temporal discretization errors interact and accumulate over time. Independent refinement of either parameter produced reduced convergence rates ($p \approx 1.06$), indicating that unresolved error from the fixed parameter influences the solution. This demonstrates that accurate error estimation in transient multiphysics simulations requires consistent co-refinement of spatial and temporal discretization. 

The interface energy consistency results support the coupled solution. The interface heat flux is more sensitive to local discretization since it does not exhibit monotonic convergence with mesh refinement. The imbalance remains small, approximately $0.01\%$ of the applied heat source, indicating that energy is well conserved across the interface. This imbalance is reasonable so long as it remains smaller than the discretization error.


% % Add this but edit cause it's too long for interface analysis:
% The absence of monotonic reduction in the normalized interface imbalance is not unexpected for the coupled conjugate heat transfer case. Unlike the isolated radial conduction verification problem, the coupled case introduces additional numerical error sources, including interface interpolation between separate fluid and solid regions, gradient reconstruction used to compute heat flux, pressure--velocity--temperature coupling, and finite solver tolerances. Because the imbalance is computed from integrated derived quantities rather than primary solution variables, it can be more sensitive to these effects and may plateau once spatial discretization error becomes smaller than the coupling and linear-solver error floor. 
% %%%%%%%%%%%%%%%%%%%%%%%%%%%%%%%

% The differing convergence behavior between the radial conduction and coupled configurations highlights the role of additional numerical error sources introduced by conjugate heat transfer coupling. In the radial conduction case, the interface imbalance is primarily governed by spatial discretization error and exhibits the expected reduction with mesh refinement. 

% In contrast, the coupled configuration introduces competing error sources, including interpolation across the fluid–solid interface, solver convergence tolerances, and coupling between the velocity and temperature fields. These errors do not scale directly with mesh refinement and can dominate the interface flux balance. Additionally, because the interface heat flux depends on temperature gradients, it is more sensitive to local discretization and interpolation errors than the primary solution variables. As a result, the interface imbalance approaches a plateau rather than exhibiting strict monotonic convergence.


The computational cost analysis further emphasizes the practical implications of verification. While refinement improves accuracy, diminishing returns are observed as computational cost increases. Importantly, the relationship between discretization parameters and computational cost is not strictly monotonic. Larger timesteps or coarser meshes do not always yield lower runtime due to increased solver iterations and numerical stiffness. This highlights the need to evaluate computational cost based on measured runtime rather than assuming a direct relationship with discretization parameters.

The selected M3–T3 configuration provides a balance between accuracy (relative error $\sim 0.1\%$) and computational cost (24 minutes), making it suitable for downstream validation studies. This choice reflects a risk-informed approach to modeling, where sufficient accuracy is achieved without exceeding practical computational constraints. If stricter error tolerances are required for a given application, further refinement can be applied; however, this will incur increased computational cost and should be justified based on the sensitivity of the decision-making context.

Geometric simplification was essential to enable rigorous verification while retaining the dominant physics of the system. Although detailed device features such as coil geometry and the expanding mouthpiece were excluded, the simplified model captures the key fluid and thermal transport behavior governing the system response. This level of abstraction allows verification to be performed under controlled conditions while maintaining relevance to the physical system. More complex geometric features can be reintroduced during validation without compromising the credibility established through verification.

Finally, this study highlights the role of reproducibility in verification. High-fidelity simulations require large datasets and substantial computational resources, particularly for fine mesh and timestep combinations. By implementing scripted workflows and providing publicly accessible computational artifacts, this work ensures that verification results are traceable, reproducible, and accessible. This is particularly important in settings where computational resources are limited, as it enables broader use and independent assessment of the modeling approach. This verification process establishes a credible foundation for subsequent validation and application of the model within the ASME V\&V 40 framework. 

This work is intentionally scoped to verification in order to isolate and quantify numerical error prior to introducing additional sources of uncertainty. While uncertainty quantification and validation are essential components of the full VVUQ process, combining all elements within a single study can obscure the attribution of error sources. By separating verification from validation and uncertainty quantification, this work establishes a controlled and transparent foundation for subsequent analysis. A companion study will extend this work to include validation and uncertainty quantification under application-relevant conditions.
